\section{Related Work}

\label{sec:related}

\paragraph{Tool-Creating Agents.}
The closest prior work is \textsc{bmo}~\cite{ngrok2024bmo}, an agent that
can generate and register new tools at runtime to address capability gaps.
\textsc{bmo} demonstrates that tool creation is practically viable but does not
address extension development, cross-component modification, or the formal safety
properties required for production deployment. Our work extends \textsc{bmo}'s
runtime creation to a full four-tier hierarchy with isolation guarantees.

\paragraph{Skill Libraries.}
Voyager~\cite{wang2023voyager} builds an ever-growing skill library for
Minecraft by having GPT-4 write JavaScript code that is verified and stored for
reuse. Each skill is tested in-game before storage. HHP generalizes this idea
beyond game-specific code to production software components, and adds the
critical missing piece: a safety framework for when the skill is not a
sandboxed game action but a modification to shared production infrastructure.

\paragraph{Automated Design of Agentic Systems.}
ADAS~\cite{hu2024adas} uses a meta-agent to search the space of agentic system
designs, producing new agent specifications expressed in code. ADAS operates
at design time, not at agent runtime. HHP addresses the complementary problem:
runtime self-modification during active task execution, where modifications must
not disrupt the running session.

\paragraph{Self-Instruct and Self-Improvement.}
Self-Instruct~\cite{wang2022self} and related work~\cite{zelikman2022star,
yuan2023scaling} demonstrate that models can generate training data and
instructions that improve themselves. These methods improve model weights;
HHP improves the harness rather than the weights, making it complementary to
and composable with weight-based self-improvement.

\paragraph{Program Synthesis and Self-Modifying Code.}
Program synthesis~\cite{gulwani2017program} generates code from
specifications. Classical self-modifying code~\cite{holland1975adaptation}
allows programs to rewrite their own instructions. Modern Genetic
Programming~\cite{koza1994genetic} evolves programs via mutation and selection.
HHP draws on these traditions but differs in that modifications are semantically
motivated (derived from observed friction), architecturally structured (four
tiers), and formally safety-bounded rather than evolutionary.

\paragraph{Git-Based Review as Safety.}
The use of pull request review as a human gate for agent-generated code
modifications has been explored in the context of automated bug
fixing~\cite{sobania2023analysis} and automated refactoring~\cite{bavishi2019autopandas}.
HHP systematizes this pattern: the PR is not merely a record but a required
gate for all Tier 4 (cross-component) modifications, with CI as an automated
co-reviewer.

\paragraph{Active Semantic Optimization.}
ASO~\cite{hanzo2026aso} and DSO~\cite{hanzo2026dso} provide the telemetry and
experience-sharing infrastructure that HHP builds on. GRPO-extracted experiential
priors in ASO encode what approaches work best for a given agent; HHP closes the
loop by allowing those insights to drive harness modifications that make the
winning approaches available as first-class tools.

% ============================================================================
% 3. THE SELF-MODIFICATION TAXONOMY
% ============================================================================
